\section{The \texorpdfstring{$(k, \ell)$}{(k, l)}-count matroid}
\label{sec:count-matroid}

Let $V$ be a finite set, let $K_V$ be the complete graph on $V$, and
let $k$ and $\ell$ be non-negative integers with $\ell < 2k$. We call a
set $I \subseteq E(K_V)$ \emph{$(k, \ell)$-sparse} if the graph on $V$
with edge set $I$ is $(k, \ell)$-sparse (\cref{def:isSparse}). In this
section we show that the $(k, \ell)$-sparse subsets of $E(K_V)$ are the
independent sets of a matroid on $E(K_V)$, the $(k, \ell)$-count
matroid. In \cref{sec:rigidity-matroid-instance} we took its case
$(k, \ell) = (2, 3)$ as the definition of the planar rigidity matroid
$\mathcal{R}(V)$. \Cref{thm:linearRigidityMatroid-eq-rigidityMatroid}
identifies $\mathcal{R}(V)$ with the rigidity matroid of a framework of
$K_V$.

The count matroid goes back to Lorea and, independently, to White and
Whiteley. Whiteley~\cite[Proposition~A.1.1]{whiteley1996} proves it in
the generality of hypergraphs and credits both. Lee and
Streinu~\cite[Theorem~2]{leeStreinu2008} prove a version for
multigraphs. Their proof rests on the structure of tight sets, which we
also use.

Tight sets are as in \cref{sec:sparsity}, where the paragraph after
\cref{def:isTightOn} relates them to Jord\'an's critical sets. Let $G$
be a $(k, \ell)$-sparse graph on $V$. Following Lee and Streinu, who
speak of the induced subgraphs, we call a $(k, \ell)$-tight set of $G$
a \emph{component} of $G$ if no other $(k, \ell)$-tight set of $G$
contains it.

A family of subsets of a finite set is the family of independent sets
of a matroid if it contains the empty set, is closed under taking
subsets, and has the augmentation property: if $I$ and $J$ are members
with $|I| < |J|$, then $I \cup \{e\}$ is a member for some
$e \in J \setminus I$. For the $(k, \ell)$-sparse subsets of $E(K_V)$
the first two properties are immediate. We prove the third in
\cref{sec:count-matroid-aug} (\cref{lem:isSparse-aug}), and define the
matroid in \cref{sec:count-matroid-def}.

We sketch the proof of the augmentation property. Let $I$ and $J$ be
$(k, \ell)$-sparse with $|I| < |J|$, and let $G$ be the graph on $V$
with edge set $I$. Since $\ell < 2k$, two distinct vertices that lie in
a tight set of $G$ lie in exactly one component of $G$. So two distinct
components share at most one vertex, and no edge lies inside both.
Suppose that $I \cup \{e\}$ is $(k, \ell)$-sparse for no
$e \in J \setminus I$. Then the endpoints of each edge of
$J \setminus I$ lie in a tight set of $G$, and hence in a component.
Inside each of these components, $J$ has at most as many edges as $I$,
since the component is tight in $G$ and $J$ is sparse. Every other edge
of $J$ is an edge of $I$. Hence $|J| \le |I|$, a contradiction.

Lee and Streinu also decide whether a graph is $(k, \ell)$-sparse by an
algorithm, the $(k, \ell)$-pebble game, which we formalize in
\cref{sec:pebble-game}.

\subsection{Components and augmentation}
\label{sec:count-matroid-aug}

The first lemma is the case $\ell < 2k$ of
\cref{thm:isTightOn-union-inter}, and it follows from Lee and Streinu's
Theorem~5(1)~\cite{leeStreinu2008}. From it we obtain the component
that contains a given set of at least two vertices. We then show that
two distinct components share at most one vertex, which is their
Theorem~5(2). These lemmas and
\cref{lem:exists-isTightOn-of-insert-not-sparse} give the augmentation
property.

\begin{lemma}[Tight sets that share two vertices]
  \label{lem:isTightOn-union-inter-of-pair}
  \lean{SimpleGraph.IsTightOn.union_inter_of_pair}
  \leanok
  \uses{def:isTightOn, thm:isTightOn-union-inter}
  Let $V$ be finite, let $\ell < 2k$, and let $G$ be a
  $(k, \ell)$-sparse graph on $V$. Let $s, t \subseteq V$ be
  $(k, \ell)$-tight in $G$, and suppose that two distinct vertices $u$
  and $v$ lie in both $s$ and $t$. Then $s \cup t$ and $s \cap t$ are
  $(k, \ell)$-tight in $G$.
\end{lemma}
\begin{proof}
  \leanok
  \uses{thm:isTightOn-union-inter}
  Since $u$ and $v$ lie in $s \cap t$, we have $|s \cap t| \ge 2$, and
  so $\ell < 2k \le k\,|s \cap t|$. Hence
  \cref{thm:isTightOn-union-inter} applies.
\end{proof}

The next lemma is one direction of Lee and Streinu's
Corollary~6~\cite{leeStreinu2008}.

\begin{lemma}[An edge that cannot be added lies in a tight set]
  \label{lem:exists-isTightOn-of-insert-not-sparse}
  \lean{SimpleGraph.IsSparse.exists_isTightOn_of_insert_not_sparse}
  \leanok
  \uses{def:isSparse, def:isTightOn}
  Let $V$ be finite, let $G$ be a $(k, \ell)$-sparse graph on $V$, and
  let $u$ and $v$ be distinct vertices that are not adjacent in $G$.
  Let $G + uv$ be obtained from $G$ by adding the edge $uv$. If
  $G + uv$ is not $(k, \ell)$-sparse, then some set containing $u$ and
  $v$ is $(k, \ell)$-tight in $G$.
\end{lemma}
\begin{proof}
  \leanok
  \uses{def:isSparse, def:isTightOn}
  Since $G + uv$ is not sparse, there is a set $S \subseteq V$ with
  $\ell \le k\,|S|$ and $|\edgesIn[G + uv]{S}| + \ell > k\,|S|$
  (\cref{def:isSparse}). If $u$ or $v$ lies outside $S$, then $G + uv$
  and $G$ have the same edges inside $S$, and this contradicts the
  sparsity of $G$ at $S$.

  So $u$ and $v$ lie in $S$. Since $uv$ is not an edge of $G$, the set
  $S$ spans exactly one more edge of $G + uv$ than of $G$. Hence $k\,|S| < |\edgesIn{S}| + 1 + \ell$, while the
  sparsity of $G$ gives $|\edgesIn{S}| + \ell \le k\,|S|$. So
  $|\edgesIn{S}| + \ell = k\,|S|$, that is, $S$ is tight in $G$.
\end{proof}

By \cref{lem:isTightOn-union-inter-of-pair}, the tight sets of a sparse
graph that contain two given vertices are closed under union. We
therefore consider the union of all tight sets that contain a given set
of vertices.

\begin{definition}[The union of the tight sets containing a set]
  \label{def:isSparse-maxBlock}
  \lean{SimpleGraph.maxBlock}
  \leanok
  \uses{def:isTightOn}
  Let $V$ be finite, let $G$ be a graph on $V$, and let
  $X \subseteq V$. We write $C_G(X)$ for the union of all sets that
  contain $X$ and are $(k, \ell)$-tight in $G$. If there is no such
  set, then $C_G(X)$ is empty.
\end{definition}

\begin{lemma}[The union of the tight sets containing a set is tight]
  \label{lem:isSparse-maxBlock-isTightOn}
  \lean{SimpleGraph.IsSparse.maxBlock_isTightOn}
  \leanok
  \uses{def:isSparse-maxBlock, lem:isTightOn-union-inter-of-pair}
  Let $V$ be finite, let $\ell < 2k$, and let $G$ be a
  $(k, \ell)$-sparse graph on $V$. Let $X \subseteq V$ have at least
  two elements, and suppose that some $(k, \ell)$-tight set of $G$
  contains $X$. Then $C_G(X)$ is $(k, \ell)$-tight in $G$.
\end{lemma}
\begin{proof}
  \leanok
  \uses{lem:isTightOn-union-inter-of-pair}
  Let $\mathcal{F}$ be the family of tight sets of $G$ that contain
  $X$. It is finite, since $V$ is, and it is non-empty by hypothesis.
  Choose distinct vertices $u$ and $v$ in $X$. Every member of
  $\mathcal{F}$ contains $u$ and $v$, so the union of two members of
  $\mathcal{F}$ is tight (\cref{lem:isTightOn-union-inter-of-pair}) and
  is again a member of $\mathcal{F}$. By induction on the number of
  members, the union of all members of $\mathcal{F}$ is a member of
  $\mathcal{F}$. This union is $C_G(X)$, which is therefore tight.
\end{proof}

Under these hypotheses, $C_G(X)$ is therefore the largest tight set of
$G$ that contains $X$. It is a
component of $G$: a tight set that contains $C_G(X)$ also contains $X$,
and so it is contained in $C_G(X)$.

\begin{lemma}[Distinct components share at most one vertex]
  \label{lem:isSparse-maxBlock-edge-disjoint}
  \lean{SimpleGraph.IsSparse.maxBlock_eq_of_subset_maxBlock}
  \leanok
  \uses{def:isSparse-maxBlock, lem:isSparse-maxBlock-isTightOn,
        lem:isTightOn-union-inter-of-pair}
  Let $V$ be finite, let $\ell < 2k$, and let $G$ be a
  $(k, \ell)$-sparse graph on $V$. Let $X, Y \subseteq V$ each have at
  least two elements. Suppose that some $(k, \ell)$-tight set of $G$
  contains $X$, and that $Y \subseteq C_G(X)$. Then
  $C_G(Y) = C_G(X)$. In particular, two distinct components of $G$
  share at most one vertex, and so no edge.
\end{lemma}
\begin{proof}
  \leanok
  \uses{lem:isTightOn-union-inter-of-pair,
        lem:isSparse-maxBlock-isTightOn}
  The set $C_G(X)$ is tight (\cref{lem:isSparse-maxBlock-isTightOn}).
  It contains $Y$, so it is one of the sets whose union is $C_G(Y)$.
  Hence $C_G(X) \subseteq C_G(Y)$. Since some tight set contains $Y$,
  the set $C_G(Y)$ is tight as well
  (\cref{lem:isSparse-maxBlock-isTightOn}). Both $C_G(X)$ and $C_G(Y)$
  contain $Y$, which has at least two elements, so their union is tight
  (\cref{lem:isTightOn-union-inter-of-pair}). This union contains $X$,
  so it is contained in $C_G(X)$. Hence $C_G(Y) \subseteq C_G(X)$.

  For the last claim, let $C$ be a component of $G$ with at least two
  vertices. Then $C \subseteq C_G(C)$, and $C_G(C)$ is tight
  (\cref{lem:isSparse-maxBlock-isTightOn}). Since no other tight set
  contains $C$, we have $C = C_G(C)$. If distinct vertices $u$ and $v$
  lie in $C$, then the first claim, with $X = C$ and $Y = \{u, v\}$,
  gives $C = C_G(\{u, v\})$. So two components that share $u$ and $v$
  are both equal to $C_G(\{u, v\})$.
\end{proof}

We can now prove the augmentation property.

\begin{lemma}[Augmentation]
  \label{lem:isSparse-aug}
  \lean{SimpleGraph.IsSparse.exists_aug_of_lt_two_mul}
  \leanok
  \uses{def:isSparse, def:isTightOn,
        lem:isTightOn-union-inter-of-pair,
        lem:exists-isTightOn-of-insert-not-sparse,
        def:isSparse-maxBlock,
        lem:isSparse-maxBlock-isTightOn,
        lem:isSparse-maxBlock-edge-disjoint}
  Let $V$ be finite, let $\ell < 2k$, and let $I, J \subseteq E(K_V)$
  be $(k, \ell)$-sparse with $|I| < |J|$. Then there is an edge
  $e \in J \setminus I$ such that $I \cup \{e\}$ is $(k, \ell)$-sparse.
\end{lemma}
\begin{proof}
  \leanok
  \uses{def:isSparse, def:isTightOn,
        lem:isTightOn-union-inter-of-pair,
        lem:exists-isTightOn-of-insert-not-sparse,
        def:isSparse-maxBlock,
        lem:isSparse-maxBlock-isTightOn,
        lem:isSparse-maxBlock-edge-disjoint}
  Let $G$ and $H$ be the graphs on $V$ with edge sets $I$ and $J$, so
  that $G$ and $H$ are $(k, \ell)$-sparse. Suppose, to the contrary,
  that $I \cup \{e\}$ is $(k, \ell)$-sparse for no
  $e \in J \setminus I$. Let $e = uv$ be an edge in $J \setminus I$.
  Since $e \notin I$, the vertices $u$ and $v$ are not adjacent in $G$,
  and $G + uv$ is the graph on $V$ with edge set $I \cup \{e\}$. It is
  not sparse, so some tight set of $G$ contains $u$ and $v$
  (\cref{lem:exists-isTightOn-of-insert-not-sparse}). Let
  $C_e = C_G(\{u, v\})$. This set contains $u$ and $v$, and it is tight
  in $G$ (\cref{lem:isSparse-maxBlock-isTightOn}).

  Let $\mathcal{C}$ be the family of the sets $C_e$, for
  $e \in J \setminus I$. An edge $xy$ of $K_V$ lies inside at most one
  member of $\mathcal{C}$. Indeed, if $x$ and $y$ lie in $C_e$, then
  $C_G(\{x, y\}) = C_e$ (\cref{lem:isSparse-maxBlock-edge-disjoint}).
  Let $I_0$ and $J_0$ be the sets of edges of $I$ and of $J$ that lie
  inside no member of $\mathcal{C}$. Then $I$ is the disjoint union of
  $I_0$ and the sets $\edgesIn{C}$, for $C \in \mathcal{C}$, and
  likewise for $J$ and $H$. So
  \[
    |I| = \sum_{C \in \mathcal{C}} |\edgesIn{C}| + |I_0|
    \quad\text{and}\quad
    |J| = \sum_{C \in \mathcal{C}} |\edgesIn[H]{C}| + |J_0|.
  \]

  We compare the two sums term by term. Every edge $e$ of
  $J \setminus I$ lies inside $C_e$. So every edge of $J_0$ is an edge
  of $I$, and $J_0 \subseteq I_0$. Let $C \in \mathcal{C}$. Since $C$ is
  tight in $G$, we have $|\edgesIn{C}| + \ell = k\,|C|$, and in
  particular $\ell \le k\,|C|$. So the sparsity of $H$ at $C$ gives
  \[
    |\edgesIn[H]{C}| + \ell \;\le\; k\,|C| \;=\; |\edgesIn{C}| + \ell.
  \]
  Summing, we get $|J| \le |I|$, which contradicts $|I| < |J|$.
\end{proof}

\subsection{The matroid}
\label{sec:count-matroid-def}

The $(k, \ell)$-sparse subsets of $E(K_V)$ satisfy the independence
axioms of a matroid. The empty set is $(k, \ell)$-sparse, since no set
of vertices spans an edge of the graph on $V$ with no edges
(\cref{lem:edgesIn-corners}). A subset of a $(k, \ell)$-sparse set is
$(k, \ell)$-sparse (\cref{lem:isSparse-mono-left}). The augmentation
property is \cref{lem:isSparse-aug}.

\begin{definition}[$(k, \ell)$-count matroid]
  \label{def:countMatroid}
  \lean{SimpleGraph.countMatroid}
  \leanok
  \uses{def:isSparse, lem:isSparse-aug, lem:edgesIn-corners,
    lem:isSparse-mono-left}
  Let $V$ be finite, and let $k$ and $\ell$ be non-negative integers
  with $\ell < 2k$. The \emph{$(k, \ell)$-count matroid} on $V$ is the
  matroid on $E(K_V)$ whose independent sets are the sets
  $I \subseteq E(K_V)$ for which the graph on $V$ with edge set $I$ is
  $(k, \ell)$-sparse.
\end{definition}
\begin{fmlnote}
  Formally, the count matroid is a matroid on the set of all unordered
  pairs of elements of $V$, the pairs $\{v, v\}$ included, with ground
  set $E(K_V)$. So a set of pairs is independent only if it is
  contained in $E(K_V)$.
\end{fmlnote}

\begin{theorem}[Independent sets of the count matroid]
  \label{thm:countMatroid-indep-iff}
  \lean{SimpleGraph.countMatroid_indep_iff}
  \leanok
  \uses{def:countMatroid, def:isSparse}
  Let $V$ be finite, let $k$ and $\ell$ be non-negative integers with
  $\ell < 2k$, and let $I$ be a set of unordered pairs of elements of
  $V$. Then $I$ is independent in the $(k, \ell)$-count matroid on $V$
  if and only if $I \subseteq E(K_V)$ and the graph on $V$ with edge
  set $I$ is $(k, \ell)$-sparse.
\end{theorem}
\begin{proof}
  \leanok
  \uses{def:countMatroid}
  The independent sets of the count matroid are defined to be these
  sets (\cref{def:countMatroid}).
\end{proof}
